\chapter{Build Pipeline}
\label{ch:pipeline}

The Veilbox build pipeline transforms a bare Linux kernel source tree and
a collection of binary tarballs into a bootable disk image. This chapter
documents every stage of the pipeline, explaining the purpose of each step,
the tools involved, and the output produced.

\section{Pipeline Overview}

The build is orchestrated by \texttt{build.sh}, a 759-line Bash script
divided into 14 sequential stages. Each stage is implemented as a
standalone function that can be called independently.

\begin{figure}[H]
\centering
\begin{tikzpicture}[node distance=0.8cm]
\node[box] (deps) {Install Deps};
\node[box, below of=deps] (dirs) {Create Rootfs};
\node[box, below of=dirs] (config) {Write Config};
\node[box, below of=config] (kconf) {Configure Kernel};
\node[box, below of=kconf] (strip) {Strip Drivers};
\node[box, below of=strip] (bins) {Populate Binaries};
\node[box, below of=bins] (ssh) {SSH Keys};
\node[box, below of=ssh] (libs) {Copy Libraries};
\node[box, below of=libs] (cpio) {Generate Initramfs};
\node[box, below of=cpio] (kernel) {Build Kernel};
\node[box, below of=kernel] (sqfs) {SquashFS};
\node[box, below of=sqfs] (state) {State Image};
\node[box, below of=state] (disk) {Bootable Disk};
\node[box, below of=disk] (vdi) {Convert VDI};
\draw[arrow] (deps) -- (dirs);
\draw[arrow] (dirs) -- (config);
\draw[arrow] (config) -- (kconf);
\draw[arrow] (kconf) -- (strip);
\draw[arrow] (strip) -- (bins);
\draw[arrow] (bins) -- (ssh);
\draw[arrow] (ssh) -- (libs);
\draw[arrow] (libs) -- (cpio);
\draw[arrow] (cpio) -- (kernel);
\draw[arrow] (kernel) -- (sqfs);
\draw[arrow] (sqfs) -- (state);
\draw[arrow] (state) -- (disk);
\draw[arrow] (disk) -- (vdi);
\end{tikzpicture}
\caption{Veilbox build pipeline — 14 stages from environment setup to
bootable VirtualBox disk image. Arrows indicate data flow.}
\label{fig:pipeline}
\end{figure}

\section{Stage 1: Install Dependencies}

The \texttt{install\_deps()} function checks for required tools and
installs them using the host package manager.

\begin{lstlisting}[language=bash]
install_deps() {
    local pkgs=(gcc make flex bison bc openssl-devel ...)
    if command -v dnf &>/dev/null; then
        sudo dnf install -y "${pkgs[@]}"
    elif command -v apt &>/dev/null; then
        sudo apt install -y "${pkgs[@]}"
    else
        echo "Warning: unsupported package manager"
        echo "Please install: ${pkgs[*]}"
    fi
}
\end{lstlisting}

\section{Stage 2: Create Rootfs Directories}

The standard Linux filesystem hierarchy is created under \texttt{rootfs/}.

\begin{lstlisting}[language=bash]
create_rootfs_dirs() {
    local dirs=(
        rootfs/bin   rootfs/sbin   rootfs/etc
        rootfs/etc/init.d   rootfs/etc/dropbear
        rootfs/usr/bin      rootfs/usr/sbin
        rootfs/lib64        rootfs/mnt/state
        rootfs/proc         rootfs/sys
        rootfs/dev          rootfs/tmp
        rootfs/var/log      rootfs/var/run
        rootfs/run          rootfs/root
    )
    mkdir -p "${dirs[@]}"
}
\end{lstlisting}

\section{Stage 3: Write Configuration Files}

This stage generates the system configuration files that define Veilbox's
behavior. The files written include:

\begin{itemize}[nosep]
  \item \texttt{/etc/inittab}: BusyBox init configuration — spawns getty
        on tty1 and ttyS0, runs rcS at boot, handles power events.
  \item \texttt{/etc/init.d/rcS}: System initialization script — mounts
        virtual filesystems, starts network, launches services.
  \item \texttt{/etc/passwd}: Root user entry with shell set to /bin/sh.
  \item \texttt{/etc/shadow}: Hashed password for root (veiladmin) using
        SHA-512 crypt hash.
  \item \texttt{/etc/hostname}: The string \texttt{veilbox}.
  \item \texttt{/etc/containerd/config.toml}: Containerd configuration
        with snapshotter, namespace, and registry settings.
  \item \texttt{/root/.profile}: Shell profile with colored prompt and
        IP address display.
  \item \texttt{/usr/share/udhcpc/default.script}: DHCP client script
        for interface configuration.
\end{itemize}

\section{Stage 4: Configure Kernel}

\subsection{Default Configuration}

The kernel configuration starts from \texttt{make defconfig}, which
generates a configuration file with approximately 10,000 options enabled
based on the target architecture.

\subsection{Config Fragment Merge}

A custom configuration fragment (\texttt{kernel/configs/custom-os.config})
is merged on top of the defconfig baseline using the kernel's
\texttt{merge\_config.sh} script:

\begin{lstlisting}[language=bash]
configure_kernel() {
    make defconfig
    ./scripts/kconfig/merge_config.sh \
        -O .config \
        .config \
        kernel/configs/custom-os.config
    make olddefconfig
}
\end{lstlisting}

The \texttt{merge\_config.sh} tool applies the fragment on top of the
existing .config, enabling options set to \texttt{y} or \texttt{m} in
the fragment and disabling options set to \texttt{n}. The
\texttt{olddefconfig} pass resolves any new dependencies introduced by
the fragment.

\section{Stage 5: Strip Unnecessary Drivers}

After merging the config fragment, a series of \texttt{sed} commands
removes kernel subsystems that are not needed in a virtualized container
host:

\begin{lstlisting}[language=bash]
strip_unnecessary_drivers() {
    # No sound support in VM
    sed -i '/CONFIG_SOUND[= ]/c\# CONFIG_SOUND is not set' .config

    # No USB needed
    sed -i '/CONFIG_USB[= ]/c\# CONFIG_USB is not set' .config

    # No Wi-Fi or Bluetooth
    sed -i '/CONFIG_WLAN[= ]/c\# CONFIG_WLAN is not set' .config

    # No ATA or SCSI drivers
    sed -i '/CONFIG_ATA[= ]/c\# CONFIG_ATA is not set' .config

    make olddefconfig  # Fix dependency breakage
}
\end{lstlisting}

\section{Stage 6: Populate Rootfs Binaries}

This stage downloads and installs the four user-space binaries that
constitute the Veilbox userspace:

\begin{enumerate}[nosep]
  \item \textbf{BusyBox} (1.35.0): \texttt{busybox} binary is copied to
        \texttt{rootfs/bin/busybox}. The installation step also creates
        the standard symlinks (\texttt{/bin/sh}, \texttt{/bin/ls}, etc.).
  \item \textbf{containerd} (1.7.13): Extracted from tar.gz, the
        \texttt{containerd}, \texttt{containerd-shim}, and
        \texttt{containerd-shim-runc-v2} binaries are placed in
        \texttt{rootfs/usr/bin/}.
  \item \textbf{runc} (1.1.12): Single binary placed in
        \texttt{rootfs/usr/bin/runc}.
  \item \textbf{nerdctl} (1.7.5): Single binary placed in
        \texttt{rootfs/usr/bin/nerdctl}.
  \item \textbf{Dropbear} (2024.85): Compiled from source with
        \texttt{--disable-zlib}, producing dropbear, dropbearkey,
        and dropbearconvert in \texttt{rootfs/usr/bin/}.
\end{enumerate}

\section{Stage 7: Generate SSH Keys}

An ED25519 key pair is generated for root SSH access:

\begin{lstlisting}[language=bash]
build_ssh_keys() {
    dropbearkey -t ed25519 -f rootfs/etc/dropbear/dropbear_ed25519_host_key
    dropbearkey -t rsa    -f rootfs/etc/dropbear/dropbear_rsa_host_key
    # Create authorized_keys for public key auth
    dropbearkey -t ed25519 -f /tmp/ssh_test_key \
        | grep "^ssh-ed25519" \
        > rootfs/etc/dropbear/authorized_keys
}
\end{lstlisting}

\section{Stage 8: Copy Shared Libraries}

BusyBox is statically linked, but containerd, runc, nerdctl, and Dropbear
are dynamically linked against glibc. This stage finds their shared library
dependencies and copies them to \texttt{rootfs/lib64/}:

\begin{lstlisting}[language=bash]
copy_libraries() {
    for bin in rootfs/usr/bin/*; do
        ldd "$bin" 2>/dev/null | while read -r line; do
            lib=$(echo "$line" | grep -oP '(?<==> ).*\.so[^ ]*' || true)
            if [ -n "$lib" ] && [ -f "$lib" ]; then
                cp "$lib" rootfs/lib64/
            fi
        done
    done
}
\end{lstlisting}

\section{Stage 9: Generate Initramfs File List}

The initramfs is specified as a kernel file list — a text file that maps
file paths to their content, permissions, and ownership. The
\texttt{gen\_initramfs\_list()} function converts the \texttt{rootfs/}
directory into this format:

\begin{lstlisting}[language=bash]
gen_initramfs_list() {
    # Use the kernel's gen_initramfs_list.sh script
    scripts/gen_initramfs_list.sh rootfs/ > initramfs.list

    # Add device nodes
    echo "nod /dev/console 0600 0 0 c 5 1"   >> initramfs.list
    echo "nod /dev/tty0   0600 0 0 c 4 0"    >> initramfs.list
    echo "nod /dev/ttyS0  0600 0 0 c 4 64"   >> initramfs.list
    echo "nod /dev/null   0666 0 0 c 1 3"    >> initramfs.list
    echo "nod /dev/zero   0666 0 0 c 1 5"    >> initramfs.list
}
\end{lstlisting}

\section{Stage 10: Build Kernel}

The kernel is compiled with \texttt{make -j\$(nproc)}. The embedded
initramfs is included automatically because the config specifies
\texttt{CONFIG\_INITRAMFS\_SOURCE="initramfs.list"}.

\begin{lstlisting}[language=bash]
build_kernel() {
    make -j$(nproc) bzImage
    cp arch/x86/boot/bzImage output/vmlinuz
}
\end{lstlisting}

\section{Stage 11: Create SquashFS Reference}

An optional compressed rootfs archive is created for inspection:

\begin{lstlisting}[language=bash]
mksquashfs rootfs output/rootfs.squashfs -comp xz
\end{lstlisting}

\section{Stage 12: Create State Image}

A 128 MB ext4 filesystem is created for persistent state storage:

\begin{lstlisting}[language=bash]
qemu-img create -f raw output/state.img 128M
mkfs.ext4 -F output/state.img
\end{lstlisting}

\section{Stage 13: Create Bootable Disk Image}

This is the most complex stage. A 8 GB raw disk image is created,
partitioned, formatted, and populated with the bootloader and kernel.
The detailed process is described in Chapter 6 (GRUB).

\section{Stage 14: Convert to VDI}

Finally, the raw image is converted to VirtualBox format:

\begin{lstlisting}[language=bash]
qemu-img convert -f raw -O vdi output/veilbox.raw output/veilbox.vdi
\end{lstlisting}


\section{Detailed Stage Analysis}

\subsection{Dependency Detection}

The build script's \texttt{install\_deps()} function uses a combination
of \texttt{command -v} checks and feature tests to determine whether each
required tool is available:

\begin{lstlisting}[language=bash]
check_deps() {
    local missing=0
    for cmd in gcc make flex bison bc curl rsync cpio xz \
               qemu-img debugfs sfdisk grub2-mkimage; do
        if ! command -v "$cmd" &>/dev/null; then
            echo "Missing: $cmd"
            missing=$((missing+1))
        fi
    done
    return $missing
}
\end{lstlisting}

\subsection{Error Handling Strategy}

The build script uses a centralized error handling pattern:

\begin{lstlisting}[language=bash]
set -euo pipefail  # Exit on error, undefined vars, pipe failures

# Each stage function returns non-zero on failure
build_kernel() {
    echo "==> Building kernel..."
    make -j$(nproc) bzImage || {
        echo "ERROR: Kernel build failed"
        echo "Check kernel config with: make menuconfig"
        return 1
    }
}
\end{lstlisting}

\subsection{Caching Strategy}

To avoid redundant downloads and rebuilds, the script checks for
existing artifacts:

\begin{lstlisting}[language=bash]
# Skip kernel extraction if already present
if [ -d linux-7.1.0-rc6 ]; then
    echo "Kernel source already extracted, skipping..."
else
    tar xf linux-7.1.0-rc6.tar.xz
fi

# Skip busybox download if already cached
if [ -f busybox-1.35.0.tar.bz2 ]; then
    echo "BusyBox source cached, skipping download..."
fi
\end{lstlisting}

\subsection{Clean Build Mode}

The \texttt{--clean} flag removes all build artifacts:

\begin{lstlisting}[language=bash]
clean_all() {
    rm -rf output/* linux-7.1.0-rc6 rootfs/* \
           initramfs.list .config .config.custom \
           busybox-1.35.0* containerd-1.7.13* \
           runc.amd64* nerdctl-1.7.5* dropbear-2024.85*
    echo "All build artifacts removed."
}
\end{lstlisting}

\section{Build Output Verification}

After each build stage, the script should verify the output:

\begin{lstlisting}[language=bash]
# Verify kernel binary exists and has valid format
verify_kernel() {
    local kernel="output/vmlinuz"
    if [ ! -f "$kernel" ]; then
        echo "ERROR: Kernel binary not found!"
        return 1
    fi
    local size=$(stat -c%s "$kernel")
    echo "Kernel: $kernel ($size bytes)"

    # Check for valid bzImage magic
    if ! file "$kernel" | grep -q "bzImage"; then
        echo "ERROR: Invalid kernel format!"
        return 1
    fi
}

# Verify disk image is bootable
verify_disk() {
    local disk="output/veilbox.raw"
    sfdisk -l "$disk" > /dev/null 2>&1 || {
        echo "ERROR: Invalid partition table!"
        return 1
    }
    # Check MBR boot signature
    local sig=$(od -A n -t x1 -j 510 -N 2 "$disk")
    if [ "$sig" != " 55 aa" ]; then
        echo "ERROR: Missing MBR boot signature!"
        return 1
    fi
}
\end{lstlisting}
